\subsection{The segment lattice and the last sampled degree}
\label{sec:segment-window}
\hypertarget{source-segment-window}{}

For the segment application, the variables in \eqref{eq:ties} range over an enlarged integer box. The final sampled degree contributes a separate rank loss.
\begin{paramdefinition}[The sampled degrees and size conditions for a segment]\label{hyp:segment-sampling}
Write \(s\) for the number of odd tops
between successive even tops. Put \(r=k-1\), \(P_i=p^{G-d_i}\)
where \(G\) is the length of this segment and \(d_i\) are the
increasing distances of its odd tops from its start. Its \(s+1\)
sampled degrees are \(a_i\) at those tops and \(a_+\) at the next
even top. They range over the same enlarged integer interval
\[
 W=[a,\lfloor a+a/(kY)+x_a\rfloor],
\]
where
\(0\le x_a\le2\mathcal L b'/k\) and
\(\mathcal L=C_0N_1^2\log a\). The \(b_i,b_+\) are held fixed
in the enlarged second interval.

In particular \(b_i,b_+\le5b'/2\).
Write \(\lambda=1+\log_p(8a+r)\), reserve \(N\) for the number of
tops in the cycle, and put \(Q=p^{G'}\), \(G\le G'\).
The local size hypotheses include
\[
 k\ge5,\; 2s+2\le r,\; s\le N_1,\;
 p\ge k^2Y^2,\quad pb'<a,\quad kQ\le b',\quad
 4\mathcal L<k,\quad kY^2\log a\le a/b'.
\]
The parameter \(a\) in the integer construction is the rounded
lower endpoint. We use the inequalities on that integer.
\end{paramdefinition}
The host realization uses the good cycle of Section~\ref{subsec:frobenius-selection}. The estimates preceding Proposition~\ref{prop:host-large-prime-test} verify these inequalities with the rounded endpoint and the enlarged box. In particular the two growing scales there satisfy
\[
        a/b'=X',\qquad \log Q=o(\log b').
\]


\begin{lemma}[One lattice prescription on a digit box]\label{lem:segment-common-digit-lattice}
Fix the second degrees and a nonempty digit class. On the
product of admissible first-degree progressions both the
retained label conditions and the allowed levels
\(\mathbf j\) are independent of the sampled first degrees.
The last shifted vector order \(o_+\) is less than \(a\).
On the lattice at each permitted level one has
\[
 B_0(u,v)=B_0(v,u)\qquad(u,v\in L_{\mathbf j}).
\]

\end{lemma}

\begin{proof}
\proofstep{The integral digit class.}
A digit class consists of vectors \(\rho_i,\tau_i\) of length \(k\)
with entries from \(0\) to \(P_i-1\). Write \(|\rho|=\sum_c\rho_c\).
Retain the class at these degrees when both coarse degrees are nonnegative integers, where
\[
 A_i=(b_i-|\rho_i|)/P_i,\qquad C_i=(a_i-|\tau_i|)/P_i .
\]
Put \(A_0^{\,*}=\lfloor(a-E_0)/p^G\rfloor\), using a star to
distinguish this starting capacity from the individual \(A_i\).
The odd tops split at an index \(\sigma\) into an
initial list and a final list of \(t=s-\sigma\) tied pairs
(\(0\le\sigma\le s\)).

For a block label use integers
\((w_i,h_i)\) at the initial positions, \(o_i\) at tied positions,
and \(o_+\) at the end. Set \(u_i=w_i+h_i\) or \(o_i\), respectively,
and \(K_l=A_0^{\,*}+\sum_{i<l}(A_i-u_i)\), numbering positions from zero.
The retained labels are the nonnegative choices satisfying
\begin{align*}
 w_i&\le A_i,\quad h_i\le C_i,\quad
 w_i+h_i\le A_i(1+4/p^2)+\mathcal L K_i/k,\quad
 h_i\le2\mathcal L A_0^{\,*}/k &&(i<\sigma),\\
 o_i&\le A_i,\quad o_i\le K_i+\mathcal L A_i/k &&(i\ge\sigma),\\
 o_+&\le K_s+(\mathcal L/k+4/p^2)\sum_i2a/P_i+\mathcal L b_+/k .
\end{align*}
For a nonempty class, the admissible \(a_i\)'s form a product of
arithmetic progressions of steps \(P_i\). On that entire product,
\[
 P_iC_i\ge a-kP_i\ge a-b',
 \qquad P_iA_i\le5b'/2,\qquad
 P_i\,2\mathcal L A_0^{\,*}/k\le a/2.
\]

Since \(pb'<a\) and \(p\ge k^2Y^2\), these imply both
\(C_i\ge A_i\) and \(C_i\ge 2\mathcal L A_0^{\,*}/k\).
Thus the restriction \(h_i\le C_i\) follows from its fixed
cutoff, and each minimum defining an internal tied level
is its fixed vector degree \(o_i\).
Also \(o_+\le a\). For this, \(A_0^{\,*}\le a/25\),
\(K_s\le a/25+5sb'/2\), \(sb'\le a/8\), and
\[
 \sum P_i^{-1}\le p/(p-1)\le25/24,
\]
by the strictly decreasing
powers. Inserting these and \(b_+\le5b'/2\), \(kb'\le a/4\),
in the displayed upper bound gives \(o_+<a\).

The chain now has \(2s+1\) alternating vector and covector letters,
starting and ending with a vector. Their degrees are
\((w_i,C_i-h_i)\) in the initial list, \((o_i,C_i)\) in the
tied list, and \(o_+\) at the end. Adjoin a covector \(\xi\) of
degree \(b_+\). The differences between the sampled and tied vector degrees give the decorations
\[
 z_i=A_i-o_i,\qquad z_+=a_+-o_+.
\]

The capacities at the prefixes ending in a tied vector and at
the end of the chain are \(K_i,K_s\). Subtract the prescribed
common margin \(m_0\) and take floors there to obtain the cutoffs
on first highest weights.

\proofstep{The level parameters and their order.}
A level consists of integers \(j_i\) for the ties, including the
last one \(j_+\), between zero and the minimum of the two
corresponding degrees. Subtract \(j_i\) from both degrees of that
pair. The index set of levels is fixed over the class since
\(C_i\ge A_i\ge o_i\) at the tied positions. Let \(V\) be the
constrained chain space at the reduced degrees and \(L_V\) its
coefficient lattice over \(\mathbb Z_p\). Before restoring the scalar
variables the lattice at a level is
\[
 L_{\mathbf j}=L_V\otimes_{\mathbb Z_p}
       \mathbb Z_p[\xi]_{b_+-j_+}.
\]
The factors use disjoint variables.

Write \(d_F\) for the product
of the factorials of all the reduced degrees, including the
degree of \(\xi\). Write \(T\) for the product of the \(t\) internal
operators \(\Phi_{z_i}\) on the reduced pairs.
The two forms required on \(L_{\mathbf j}\) are
\[
 B_0(f,g)=\langle Tf,g\rangle/d_F,\qquad
 B_z(f,g)=\langle\Phi_z Tf,g\rangle/d_F,\quad z=a_+-o_+ .
\]
The first form is symmetric since every \(U^lD^l\) is self-adjoint
and the operators on distinct ties commute.

\end{proof}

\begin{proposition}[Uniform mean on a segment]\label{prop:segment-discard}
Fix a segment \(I\) with the degree and size conditions of
Definition~\ref{hyp:segment-sampling}. Hold fixed the margin of
the lattice prescription in Lemma~\ref{lem:segment-common-digit-lattice}
and the \(s+1\) second weights. With \(s+1\) sampled first-weight
coordinates, the equal-weight average of the discarded
fractions over the enlarged interval box satisfies
\[
 \mathbb E_{\mathbf a,a_+}\kappa_I\le C_A(s+1)N_1/Y
\]
at the depth \eqref{eq:segment-mean-depth}. The same absolute constant
works for the depth allowance and the mean.
\end{proposition}

The proof uses the several-pair estimate separately at the
internal degrees and at the last degree. It then transports
the resulting ranks to the factorial lattice and changes back
from the auxiliary sampling law to equal weights on the box.

\begin{proof}
\proofstep{The internal degrees.}
Freeze the untied coordinates \(a_i\) and the digits and level.
At a tied position with radix \(P=P_i\) the new covector degree
\(n_i=(a_i-|\tau_i|)/P-j_i\) runs through an interval of
consecutive integers of length \(H_i\). Its law is uniform there.
One has \(H_iP\ge a/(2kY)\), since the interval in \(a_i\)
has length at least \(a/(kY)-1\), and
\(2(1+P)kY\le4b'Y\le a\). All points of this residue class above
\(a\) exceed \(|\tau_i|\).

Let \(S_i^+\) be the prefix degree
through its vector at the largest values of the other tied
coordinates. Then
\[
 P S_i^+\le 5b'/2+(2s+1)2a/p,\qquad
 P\min n_i\ge a-kP-5b'/2.
\]
The first bound uses \(P_i/P_l\le1/p\) for every earlier position
and the bound \(2a\) on each degree before division by its radix.
The vector at \(i\) has \(Po_i\le b_i\).

The second bound follows
from \(Pj_i\le b_i\) and \(|\tau_i|<kP\).
They imply \(S_i^+\le\min n_i\), since \(kP\le b'\),
\(b'<a/p\), \(2s+1\le k-2\) and \(p\ge k^2\).
They also give
\[
 S_i^+/H_i\le7/Y
\]
by multiplying the
first upper bound by \(2kY^2/a\) and using \(p\ge k^2Y^2\).
All the reduced degrees and internal decorations are at most \(2a\).
The last vector satisfies the stronger bound \(o_+\le a\) above.
The degree-size condition of \eqref{eq:ties} follows from
\(2s+2\le r\) and \(p\ge k^2\).

\proofstep{The allowance for the internal degrees.}
Use \(\kappa=3(s+2)r^2\lambda\), \(E=Y\kappa\), and write
\begin{equation}\label{eq:segment-mean-depth}
 d=t(Y+1)\kappa+\kappa,\qquad d'=\lfloor Yd\rfloor.
\end{equation}
In \eqref{eq:ties} the required lower bound on \(\kappa\) for
\(2s+1\) letters is \(\tfrac32(2s+3)r^2\lambda\).
The coefficient of each endpoint dimension is at most \(31/Y\)
because its three terms are bounded by \(14/Y,2/Y\) and
\(14\lambda/(Y^2\kappa)\).

Write \(b_V\) for the failed rank on
\(L_V\) at depth \(\lfloor d\rfloor\) and \(\widehat a_i\) for
replacement of \(a_i\) alone by its largest compatible value.
Averaging first over these progressions, then over the frozen
coordinates, gives
\begin{equation}\label{eq:seg-interior}
 \sum_{\mathbf a} b_V(\mathbf a)
   \le \frac{31}{Y}\sum_{i\ {\rm tied}}\sum_{\mathbf a}
                     \dim V(\widehat a_i).
\end{equation}
This includes \(t=0\), where the unmodified form already has
sufficient depth. Tensoring each pair of successful vectors with
each monomial \(\xi^\nu\) preserves the depth for \(B_0\).
On the partner put the integral multiplier
\((b_+-j_+)!/\prod\nu_c!\). Thus the failed rank for \(B_0\)
is at most \(b_V d_\xi\), with
\(d_\xi=\binom{b_+-j_++r-1}{r-1}\).

\proofstep{The last parameter.}
Choose an orthogonal basis of \(L_{\mathbf j}\) for \(B_0\) over
\(\mathbb Z_p\). After clearing a common denominator, a pairing
of least valuation can be realized on a single vector. If all
diagonals have greater valuation and \(B_0(u,v)\) is minimal,
use
\[
 B_0(u+v,u+v)=B_0(u,u)+2B_0(u,v)+B_0(v,v).
\]
Since \(2\) is a unit, the middle summand has strictly smallest valuation. The resulting
primitive vector splits off by subtracting integral multiples.
Induction, with a basis of the radical appended, gives the
assertion. Take the part \(f_1,\ldots,f_h\) of this basis with
nonzero diagonal of valuation at most \(\lfloor d\rfloor\).
Smith reduction identifies \(h\) with the successful rank.

Let \(o=o_+-j_+\), \(n=b_+-j_+\), \(J=\min(o,n)\),
\(q=o+n+r-1\). For the last pair write \(U=x\cdot\xi\) acting
by multiplication, \(D=\sum_c\partial_{x_c}\partial_{\xi_c}\).
Set
\[
 M_l(u,v)=
 \frac{\langle U^lD^l Tu,v\rangle}{(l!)^2\binom Jl\,d_F}\quad(0\le l\le J).
\]
The matrix with entries
\[
 p^{x_0}\sum_{l=0}^J(-1)^l
       \binom{q+z-l}{J-l} M_l(f_i,f_j)
   =p^{x_0}\binom{q+z}{J}B_z(f_i,f_j)
\]
is a polynomial in \(z\) with integral binomial-basis
coefficients for the following choice. Put \(a_*=(3s+2)a\),
\(L_*=\lfloor\log_p(8a_*+r)\rfloor+1\) and
\[
 x_0=\big((t+1)(r+1)+(2s+2)r+2(t+1)\big)L_* .
\]
For the coefficient bound, expand \(D^l\) by multi-indices
\(|\alpha|=l\). After the two divided derivatives in every
coordinate, division by \((l!)^2\) costs the inverse multinomial
\(\binom l\alpha^{-1}\). Its valuation loss is at most
\((r-1)L_*\) by the factorial formula.

The additional reciprocal
binomial costs at most \(L_*\). The same reasoning applies to
every operator in \(T\) with its own denominator
\[
 \binom{o_i'+n_i'+r+z_i-1}{l},
\]
where primes indicate
reduced degrees. Finally the monomial norm divided by \(d_F\)
costs at most one multinomial per letter, altogether at most
\(((2s+2)r+2(t+1))L_*\). Multiplication by \(U^l\) has integral
coefficients. Vandermonde's identity for each binomial in
\(q+z-l\) now proves integrality. The coefficient of \(\binom zJ\)
is exactly \(p^{x_0}B_0(f_i,f_j)\). Its diagonal entries are nonzero. Thus the coefficient of the highest power of \(z\) in the determinant is
\[
 (J!)^{-h}\prod_{i=1}^h p^{x_0} B_0(f_i,f_i).
\]
Consequently the determinant has degree \(hJ\), and the
leading coefficient has valuation at most
\(h(x_0+\lfloor d\rfloor-v_p(J!))\).

\proofstep{A Smith estimate over the last interval.}
Apply \eqref{eq:smith-average} with \(H=|W|\), shift \(a-o_+\),
\(e=\lceil\log_p(a+1)\rceil\) in place of \(h\) there, and depth
\(d'+1\). Both \(J+1,H\le a+1\).
The integral nonzero factor \(p^{x_0}\binom{q+z}{J}\) only
increases Smith valuations. Integer valuations mean that the
failed directions of the rescaled matrix at depth \(d'\) are
precisely those counted at \(d'+1\).
Legendre's formula gives
\[
 J/(p-1)-v_p(J!)\le L_*.
\]

Since \(3s+2\le p\), \(L_*\le2\lambda\), and for \(r\ge4,t\le s\)
direct multiplication gives \(x_0+L_*\le\kappa\).
Also \[
 \begin{gathered}
 J/H\le(5b'/2)kY/a\le1/(2Y), \\ e\le\kappa, \\ d\ge\kappa, \\ d'+1\ge Yd.
\end{gathered}
\]
The mean number lost on this orthogonal subsystem is therefore
at most
\[
 h\left(\frac{J}{H}+
   \frac{x_0+\lfloor d\rfloor+L_*+Je/H}{d'+1}\right)
 \le \frac{5h}{Y}.
\]
Combining with the directions already missing in \(B_0\)
bounds the mean failed rank of \(B_z\) on the full level by
\(d_\xi b_V+5d_\xi\dim V/Y\). This last average uses \(a_+\)
with \(\mathbf a,b_i,b_+\) still fixed.

\proofstep{From the levels to the factorial lattice.}
Introduce one extra pair of scalar variables \(Z_i,S_i\) at each
tie, including the last pair. In the following products \(i\)
runs over all \(t+1\) pairs. Write \(E_i=Z_iS_i+U_i\).
The rational block space is spanned by
\[
 \left(\prod_i E_i^{j_i}Z_i^{z_i}\right) f,\qquad
               f\in L_{\mathbf j},\quad \mathbf j\text{ admissible}.
\]
Extension of scalars to the coefficient field \(K=\operatorname{Frac}O\)
gives precisely the span in the definition of the segment lattice.
In fact a highest-weight cutoff is the largest prefix-stable
subspace of a monomial support condition. In fixed multidegree it is
the common kernel of maps obtained by applying words in the matrix
units and projecting to the forbidden weight spaces. There are only
finitely many independent equations, all with rational coefficients.
Kernels, their finite intersections and their dimensions commute
with field extension. The integral coefficient intersections span
these spaces over the fraction fields by clearing denominators.

\proofstep{The normalized form on one pair of letters.}
For one pair with reduced degrees \(o,n\) and decoration \(z\),
let \(H\) be harmonic projection in \(r+1\) variables at degrees
\(o+z,n\). Write \(F=D+\partial_Z\partial_S\).
For \(f\) independent of \(Z,S\),
\[
 H(Z^z f)=\sum_{l=0}^{\min(o,n)}
  \frac{(-1)^l E^l Z^z D^l f}{(l!)^2\binom{o+n+r+z-1}{l}}.
\]
Indeed \[
 F(E^l v)=E^lFv+l(r+1+\deg_x v+\deg_y v+l-1) E^{l-1}v\,;
\] the coefficients telescope and the map is the
identity modulo \(E\). The Fock adjoint of \(E\) is \(F\), so
\(H\) is self-adjoint. Replace \(Z_i^{z_i}\) in the level map
by \(f\mapsto H_i(Z_i^{z_i} f)\), taking products of these
commuting maps, followed by \(\prod E_i^{j_i}\).
This gives a map \(\Theta_{\mathbf j}\). The terms other than the
original generator have larger total level. They belong to the
indicated higher-level spaces. For a prefix including both letters
the contraction \(D_i\) is equivariant. Before the pair it acts on disjoint variables.

In either case, for the corresponding prefix action \(g\),
\[
 D_i(gf)=g(D_i f).
\]
At the intermediate prefix
it differentiates the vector. The span of the derivatives there
is an image of the old representation tensored with the dual
standard representation, so each highest weight is obtained by
subtracting one box by the dual Pieri rule. This preserves the
upper cutoff. At the last pair the covector derivative acts only
on \(\xi\). Thus the replacement is triangular, with identity
diagonal, among the rational generators. The original level spaces
are independent because in \(E_i^{j_i}\) the term of greatest \(S_i\)-degree
is \((Z_i S_i)^{j_i}\). Extract successive lexicographically
largest scalar multidegrees in a relation. Their coefficients are
exactly the separate polynomials \(f\). Summing the dimensions of these independent level spaces gives
\[
 \dim(\text{block})=\sum_{\mathbf j}d_\xi\dim V.
\]


\proofstep{Orthogonality of distinct levels.}
Images of distinct levels under \(\Theta\) are orthogonal. In a
pair with unequal levels move \(E\)'s across the pairing as
\(F\)'s until a positive power of \(F\) hits a harmonic vector.
For equal levels, moving all \(E\)'s across gives
\(j!(r+1+o+n+z)^{\overline{j}}\) at this pair. Self-adjointness
reduces the pairing of the two harmonic projections to that of
\(Z^zf\) with \(H(Z^zg)\). Only the term of \(S\)-degree zero
survives, yielding \(z!\langle f,\Phi_z g\rangle\).
Multiplying the identities gives the raw Fock pairing of the images
as \(\ell_{\mathbf j}d_F B_z\), where
\[
 \ell_{\mathbf j}=\prod_i (j_i!)^2
       \binom{o_i'+n_i'+r+z_i+j_i}{j_i} z_i!.
\]
Here \(o_i',n_i'\) are the reduced degrees of that pair.
The same operator bound used for \(M_l\) shows that
\(p^U\Theta_{\mathbf j}f\) has integral coefficients for every
\(f\in L_{\mathbf j}\), with
\[
 U=(t+1)\lceil (r+1)(1+\log_p(8a_*+r))\rceil.
\]

For this bound expand \(D_i^l\) before multiplying by the
extended \(E_i^l\), which is integral. All the relevant
factorial arguments are at most \(8a_*+r\).
This argument only needs membership in the rational block span
and the asserted coefficient bound.

\proofstep{Uniform denominator and depth bounds.}
Let \(M=\max_{\mathbf j}v_p(\ell_{\mathbf j}d_F)\).
For each of the good systems in \(B_z\) at depth \(d'\), apply
\(p^U\Theta_{\mathbf j}\) on both sides and multiply the partners
by \(p^M/(\ell_{\mathbf j}d_F)\). The scalar in the raw pairing is
\[
 p^{2U}\frac{p^M}{\ell_{\mathbf j}d_F}\,\ell_{\mathbf j}d_F\,p^{d'}=p^e,\qquad e=d'+2U+M.
\]
Orthogonality between distinct levels makes the combined matrix \(p^eI\). The factorial metric on the actual
block absorbs almost all of \(M\). More explicitly, with
\[
 v=v_p,\quad v(w)=1/2,\quad
 b_*=\prod_{i<\sigma}(A_i-w_i)!\,h_i!,
\]
that lattice consists of polynomials in the rational block span
whose monomials \(c_\nu X^\nu\) satisfy
\[
 2v(c_\nu)+v_p(\nu!)+v_p(b_*)\ge0;
\]

\(\nu!=\prod_j\nu_j!\). Its pairing is \(b_*\) times raw Fock.
For any supported \(\nu\) in this block and \emph{any} level,
the total of the factorial arguments in
\(\prod_i(j_i!)^2 z_i!\, d_F\) is exactly \(|\nu|\).
Each restored power of \(E_i\) contributes two degrees and each \(Z_i\) contributes one. Thus
\[
 |\nu|=\sum_i (2j_i+z_i)+\text{total reduced degree},
\]
the sum of the factorial arguments just listed. Legendre's formula shows at
each power of \(p\) that the difference of the sums of the
floors in these factorials and in \(\nu!\) is at most the number
of variables in \(\nu\). Each of the \(t+1\) binomials in
\(\ell_{\mathbf j}\) adds at most one at that power.

Each single
coordinate exponent is at most \(8a_*+r\), by the pairwise
bidegrees and disjoint support. Thus, writing
\(L_0=\lfloor\log_p(8a_*+r)\rfloor\),
\[
 M-v_p(\nu!)\le [(2s+2)r+3(t+1)]L_0.
\]
Since \(3s+2\le p\), \(L_0\le\lambda\) and
\(U\le(t+1)((r+1)(\lambda+1)+1)\).
It follows that
\[
 2U+M-v_p(\nu!)\le
 L_{\rm pass}=10(s+1)(r+1)\lambda+\kappa_0,\qquad
 \kappa_0=2Nk\log_p a .
\]
For \(F_0=\lfloor d'+L_{\rm pass}\rfloor\), multiply both systems
by \(w^{F_0-e-v_p(b_*)}\). The exponent can be negative.
The monomial inequality just proved, and the integrality before
this multiplication, put every resulting vector in the factorial
lattice.

Since \(v(w)=1/2\), the new cross matrix is scalar with valuation
\[
 v_p(b_*)+e+2v(w)\bigl(F_0-e-v_p(b_*)\bigr)=F_0.
\]
Smith reduction in the full lattice now gives at least
as many good directions as this scalar matrix has rows. For
example, apply Smith reduction to the full form after clearing
denominators. A restriction by integral columns and rows cannot
increase the number of factors dividing a fixed scalar. Modulo
one more uniformizer power than that scalar, the image of the
restriction factors through a module generated by the low Smith
columns. Nakayama's lemma gives the assertion. One can then take
the saturated span of the corresponding left Smith vectors in the
full lattice. Thus passage to saturation preserves the partner
criterion even if the initially constructed systems span a
nonsaturated module.

The depth needed in the final segment test is at most
\[
 Yd+L_{\rm pass}\le C_A Y^2 N_1^2 k^2\lambda+2\kappa_0=:D_I .
\]
For the inequality use \(t\le s\le N_1\), \(N_1,Y\ge1\),
\(s+2\le3N_1\), and \(d=(tY+t+1)\kappa\). A constant
\(C_A\ge56\) suffices for this depth estimate.

\proofstep{Returning to equal weights on the box.}

Write \(D_k(u)=\binom{u+k-1}{k-1}\) for \(u\ge0\), zero
otherwise. For a compatible digit class \(\delta\), put
\[
 m_\delta(\mathbf a)=D_k(A_0^{\,*})D_k(b_+)
                  \prod_i D_k(A_i)D_k(C_i).
\]
The ordinary retained-rank comparison needed here is
\begin{equation}\label{eq:ordinary-class}
 \tfrac12 m_\delta(\mathbf a)\ \le\
   \sum_{\beta\ {\rm retained\ in}\ \delta}
       \sum_{\mathbf j} d_\xi\dim V_{\beta,\mathbf j}(\mathbf a)
   \ \le\ m_\delta(\mathbf a).
\end{equation}
This comparison concerns dimensions in characteristic zero.
Its role and its required margin are recorded below.
Let \(g_\delta(\mathbf a,a_+)\) be the sum over these levels
of their failed ranks in \(B_z\) at depth \(d'\).
The preceding construction supplies a bound by \(g_\delta\) on
the failed rank of the actual block sum.

For each class the
level set is fixed over its box of compatible degrees.
Summing \eqref{eq:seg-interior} and the last-parameter estimate
therefore gives
\[
 \sum_{\mathbf a,a_+}g_\delta
 \le \frac{|W|}{Y}
      \sum_{\mathbf a}\left[31\sum_{i\ {\rm tied}}
         m_\delta(\widehat a_i)+5m_\delta(\mathbf a)\right].
\]
Only one \(D_k(C_i)\) changes at the endpoint. If it changes from
\(D_k(u)\) to \(D_k(v)\), then \(P_i u\ge a-b'\ge3a/4\),
\(P_i(v-u)\le a/(kY)+x_a\), and
\(r x_a\le 2\mathcal L b'\le a/8\).
The binomial product gives \(D_k(v)/D_k(u)\le
\exp(r(v-u)/u)\le e^3\).

\proofstep{Returning to ordinary rank weights.}
Sum over all classes and set \(m(\mathbf a)=\sum_\delta
m_\delta(\mathbf a)\). For each letter every degree-\(a_i\)
monomial has a unique expression for its exponent vector as
\(\tau_i+P_i\gamma\). Hence
\(\sum_{\tau_i}D_k(C_i)=D_k(a_i)\), and likewise for \(A_i,b_i\).
It follows that
\[
 m(\mathbf a)=D_k(A_0^{\,*})D_k(b_+)
                  \prod_i D_k(b_i)D_k(a_i).
\]
If \(C_A(t+1)N_1/Y\ge1\) the segment bound is automatic.
Otherwise \(s\le N_1\le Y\). Throughout the box
\(\max m/\min m\le\exp(sr(1/(kY)+x_a/a))\le e^3\)
when \(m>0\). Indeed \(s\le Y\),
\(rx_a\le2\mathcal L b'\), and
\(kY^2\log a\le a/b'\), \(4\mathcal L<k\) give
\(sr x_a/a\le1\) for \(a\ge e\).

The size assumptions imply
\(a\ge p b'\ge25\).
If \(D_k(A_0^{\,*})=0\), \eqref{eq:ordinary-class} forces all
ranks to vanish and the bad share is zero by definition.
Otherwise divide the total failed rank at each point by its
actual retained rank, which is at least \(m(\mathbf a)/2\).
The result for the \emph{unweighted} box average is
\begin{equation}\label{eq:operational-seg}
 \mathbb E_{\mathbf a,a_+}\kappa_I
    \le \frac{2 e^3(31 e^3 t+5)}{Y}
    \le \frac{C_A(s+1)N_1}{Y}
\end{equation}
for one absolute \(C_A\) in both the depth and the mean.
The margin and second weights were held fixed throughout. For
each segment its odd positions and its next even position are
distinct, hence each assignment of these \(s+1\) weights has the
same number of extensions to the full enlarged box. This proves
the required sampling law on the cycle. Section~\ref{sec:retained-rank} proves the ordinary-rank comparison
used here.

  The factorial-lattice calculation and the digit-class sum give the closed bound for the Gram estimate.
For the mean bound, use \(m_0=C_m kN_1Qc+12c\) and \(D_I\) as above. In its interval
\(a=\lceil a_{\rm real}\rceil\) and
\[
 x_a=\lfloor a_{\rm real}+a_{\rm real}/(kY)\rfloor+ \lfloor J_0\rfloor+c_{\rm cl}-(a+a/(kY)),
\]

where \(J_0=\mathcal L b'/k\) uses the integer lower endpoint.
Thus \(W\) is precisely the enlarged interval used in the
edge comparison. Its numerical hypotheses require
\(0\le x_a\le2J_0\), \(0\le m_0\le b'/(kQY)\).
\end{proof}
